% =====================================================================
%  D6_equilibrium_foundations.tex   (ATTEMPT 4 -- supersedes attempt 3)
%
%  Rigorous replacement for the Equilibrium foundations of
%  "Liquidity, Activism Disclosure, and Takeover Premia."
%  This file is an \input-able fragment; all macros are guarded with
%  \providecommand and the theorem environments {assumption},{remark} are
%  assumed supplied by the host (draft_v2.tex defines theorem/proposition/
%  lemma/definition; a one-line \newtheorem for assumption/remark is added at
%  integration). It uses ONLY draft-native \ref (no cleveref: the draft does
%  not load it and it is not installed in this TeX tree).
%
%  WHAT THIS FILE SUPERSEDES IN THE DRAFT:
%    - OLD Prop "Existence of Monotone Equilibrium" (self-map left unproved);
%    - OLD "Endpoint Behavior of Minority Gains" Lemma (voice collapse /
%      posteriors -> prior at kappa=1; FALSE);
%    - Assumption A7 (price injectivity), now DROPPED (bidder conditions on P);
%    - Assumption A5, split into A5a (inner price) and A5b (cross-cutoff);
%    - Assumption A6 (contraction), demoted to a Numerical Regularity;
%    - the "self-confirming" justification of the monotone-cutoff restriction.
%
%  WHAT ATTEMPT 4 FIXES RELATIVE TO ATTEMPT 3 (the refuted gaps):
%   (F1) The ill-defined measure-preserving involution iota is DELETED.
%        Endpoint symmetry is now ONLY a kappa->0,1 LIMIT statement
%        (Lemma endpoint-info + Prop endpoint-symmetry): both limits give the
%        SAME two-point D=0 law {0,pibar}; no 1-atom<->2-wing bijection, no
%        general-kappa "eq:commute", no forced reflection about 1/2.
%   (F2) NO pointwise Dmin(kappa)=Dmin(1-kappa) claim. The interior is
%        asymmetric (peak 0.59, E[pi|D=0] drift); the proof is confined to the
%        two endpoint limits.
%   (F3) The Lambda2->0-"degenerate-to-single-atom" claim is DELETED: both
%        atoms {0,pibar} persist; what vanishes near the ends is the WEIGHT.
%   (F4) The rejected midpoint chord diagnostic h(0)-2h(pibar/2)+h(pibar) is
%        NOT used. The hump engine is DERIVED by an exact decomposition
%        eq:D0decomp of the nondisclosed contribution and VERIFIED with the
%        model's own equilibrium prices.
%   (F5) SIGN CORRECTION: the realized-weight Jensen/covariance gap
%        Lambda2=(mt-m0)*P(D=0)*Cov0(pi,p) is NEGATIVE everywhere (pi deters
%        bids => Cov0(pi,p)<0 => h=pi*p convex on the realized support). The
%        gap FIGHTS the hump; it is dominated by the rising baseline term. This
%        refutes BOTH the rigor memo's "g concave" story AND the attempt-3
%        "Lambda2<0 drives the hump / U-shaped variance" story.
%   (F6) A5a stated as a MAINTAINED ASSUMPTION; the displayed 0.9474 bounds only
%        the feedback term, and with the direct term delta*p* (p*=0.847) the
%        sufficient rule exceeds 1, so the number does NOT establish |Psi'|<1.
%   (F7) Both fact sheets cited: across-equilibria multistart-at-kappa* lives in
%        phase0_facts.md; the hump/trough map + interior amplitude in
%        phase0_robustness.md.
%
%  SELF-VERIFICATION (repo's OWN solver AND equilibrium prices; scipy venv;
%   scripts _d6_check.py, _d6_engine.py, _d6_decomp.py, _d6_a5 in-line below):
%   - valid converged peak kappa*=0.59, Dmin(k*)=0.077558, amplitude 13.9% of
%     level over the converged grid (kappa=0.02 carries residual 8e-2 and is an
%     invalid near-endpoint, excluded as a maximizer).
%   - endpoint gap Dmin(0.02)=0.066778 vs Dmin(0.98)=0.066907 = 1.28e-4 on the
%     1e-5-gated solver (_d6_check.py); exact-endpoint solves agree to 6.86e-7
%     (phase0_facts.md, S1).
%   - ENGINE (computed with model prices, _d6_engine.py): mean nondisclosed bid
%     prob E0[p] and the baseline term Lam0=m0*P(D=0)*E0[p] are BOTH hump-shaped,
%     argmax at 0.59 (E0[p]: .359/.366/.404/.427/.436@.59/.421/.383/.341;
%     Lam0: .0359/.0362/.0396/.0413/.0419@.59/.0409/.0378/.0338). The gap
%     Lam2 is negative throughout (-0.0029..-0.0169 in Dmin units). Var0(pi)
%     is U-shaped (min .076 at .59), NOT the source of the hump. Both branches
%     hump: nondisclosed D0 .0599->.0774@.59->.0669 and disclosed D1 (tracking
%     omega_P) up to 1.46e-4 at .59; D0+D1=Dmin exactly (_d6_decomp.py).
%   - A5a: baseline disclosed p*=0.8473 (T=-1.025), delta*p*=0.805; with the
%     0.9474 feedback bound the sufficient rule SUMS to >1 => A5a maintained.
% =====================================================================

% --- Local macros (\providecommand: safe inside draft_v2.tex, which defines
%     \E \PP \1 \Var \Cov \R; the rest extend the draft's notation). ---
\providecommand{\E}{\mathbb{E}}
\providecommand{\PP}{\mathbb{P}}
\providecommand{\Var}{\mathrm{Var}}
\providecommand{\Cov}{\mathrm{Cov}}
\providecommand{\R}{\mathbb{R}}
\providecommand{\1}{\mathbf{1}}
\providecommand{\kp}{\kappa}
\providecommand{\kvec}{\mathbf{k}}
\providecommand{\kone}{k_1}
\providecommand{\kzero}{k_0}
\providecommand{\kD}{k_D}
\providecommand{\Tmap}{T}
\providecommand{\Psifp}{\Psi}
\providecommand{\thetastar}{\boldsymbol{\theta}^{\star}}
\providecommand{\Dmin}{\Delta^{\min}}
\providecommand{\Dact}{\Delta^{\mathrm{act}}}
\providecommand{\mt}{\tilde{m}}
\providecommand{\mR}{m^{R}}
\providecommand{\mbar}{\bar{m}}
\providecommand{\vhat}{\hat{v}}
\providecommand{\Vhat}{\hat{V}}
\providecommand{\Sbar}{\bar{S}}
\providecommand{\Kcost}{K}
\providecommand{\sigxi}{\sigma_\xi}
\providecommand{\sv}{\sigma_v}
\providecommand{\ss}{\sigma_s}
\providecommand{\deltaliq}{\delta}
\providecommand{\Ccost}{C}
\providecommand{\chivar}{\chi}
\providecommand{\pibar}{\bar{\pi}}
\providecommand{\omegaH}{\omega_H}
\providecommand{\omegaQ}{\omega_Q}
\providecommand{\omegaP}{\omega_P}
\providecommand{\xibid}{\xi}
\providecommand{\Lam}{\Lambda}

\section{Equilibrium foundations}
\label{sec:eqm-foundations}

This section gives the analytical foundation for the equilibrium objects used
throughout the paper. We (a) prove existence by exhibiting a Brouwer self-map on
the cutoff polytope; (b) establish that public voice dominates quiet voice in its
sensitivity to the signal, ordering the indifference cutoffs; (c) prove a
\emph{global} monotone--best-response lemma that delivers a cutoff partition for
\emph{any} market and bidder beliefs, replacing the earlier self-confirming
argument; (d) dispense with the auxiliary assumption that the bidder observes the
stake by conditioning the bidder on the equilibrium price directly; (e) decompose
the price/cutoff contraction hypothesis into two transparent pieces and report
the baseline margins honestly; and (f) treat uniqueness as a \emph{numerical
regularity} at the calibration rather than a theorem, while showing the
comparative statics hold across all equilibria the solver locates. We close with
endpoint symmetry of the headline object $\Dmin(\kp)$ (a $\kp\to0,1$ limit
statement) and a \emph{conditional} characterization of the interior hump whose
engine we derive and verify numerically.

% ---------------------------------------------------------------------
\subsection{The cutoff polytope and the best-response map}
% ---------------------------------------------------------------------

A pure strategy for the blockholder is a partition of the signal line into the
four action regions by an ordered triple of cutoffs. Fix a finite bracket
$[\underline{s},\overline{s}]$ containing all indifference points (its existence
is established in \ref{lem:bounded}) and define the \emph{cutoff polytope}
\begin{equation}
\label{eq:theta}
\Theta \;=\; \bigl\{\, \kvec=(\kone,\kzero,\kD)\in[\underline{s},\overline{s}]^3
\;:\; \underline{s}\le \kone \le \kzero \le \kD \le \overline{s}\,\bigr\}.
\end{equation}
$\Theta$ is the intersection of a cube with the half-spaces
$\{\kone\le\kzero\}$ and $\{\kzero\le\kD\}$; it is therefore a compact, convex,
nonempty polytope. The boundary faces $\{\kone=\kzero\}$ and $\{\kzero=\kD\}$ are
the \emph{collapse faces} on which Hold (respectively Quiet Voice) carries zero
probability; the baseline equilibrium lives on the face $\kone=\kzero$ (Hold
collapses).

The best-response map $\Tmap:\Theta\to\R^3$ sends a conjecture $\kvec$ to the
indifference cutoffs implied by it. Given $\kvec$, the market maker's pricing
function $P(\cdot,\cdot)$ and the bidder's takeover probability $p(\cdot)$ are
determined; the blockholder's payoff to action $(q,a)$ at signal $s$ is, from the
model,
\begin{equation}
\label{eq:payoff-affine}
U(q,a;s) \;=\; A(q,a)\;+\; B(q,a)\,\vhat(s)\;-\;a\,\Ccost(s),
\end{equation}
where, writing $\beta=\sv^2/\ss^2\in(0,1]$ for the posterior-mean slope and
absorbing the conjecture into the intercepts $A(\cdot)$,
\begin{equation}
\label{eq:slopes}
B(E)\;=\;0,\qquad B(H)\;=\;1,\qquad B(Q)\;=\;1,\qquad B(P)\;=\;1+\eta,
\quad \eta>0,
\end{equation}
the extra slope $\eta$ on Public Voice coming from the trading term
$q[\Vhat-P]$ with $q=+1$ (the price-taking blockholder takes $P$ as given, so the
trading channel contributes a slope of $+1$ in $\vhat$) and from the takeover
term, which is nondecreasing in $\vhat$ (\ref{lem:BP-BQ}). The component
$\Tmap_1$ ($\Tmap_2$, $\Tmap_3$) is the Exit/Hold (Hold/Quiet, Quiet/Public)
indifference signal.

\begin{lemma}[Boundedness of best responses]
\label{lem:bounded}
Under Assumptions \ref{ass:A1}, \ref{ass:A2}, \ref{ass:A4}, there is a finite
bracket $[\underline{s},\overline{s}]$, independent of the conjecture
$\kvec\in\Theta$, such that every indifference signal of $\Tmap$ lies in
$[\underline{s},\overline{s}]$.
\end{lemma}

\begin{proof}
Each indifference condition equates two payoffs of the form
\eqref{eq:payoff-affine}. The difference between two such payoffs is
\begin{equation}
\label{eq:diff}
D_{ab}(s)=\bigl[A(a)-A(b)\bigr]+\bigl[B(a)-B(b)\bigr]\vhat(s)
-\bigl[\1\{a\in\{Q,P\}\}-\1\{b\in\{Q,P\}\}\bigr]\Ccost(s).
\end{equation}
The intercepts $A(\cdot)$ are bounded uniformly over $\kvec\in\Theta$: prices lie
in $[\,0,\deltaliq\,\overline{Y}\,]$ with $\overline{Y}$ the finite maximal
terminal value, and the takeover probability $p\in[0,1]$. The slope gap
$B(a)-B(b)$ is one of $\{\pm1,\pm\eta,\pm(1+\eta)\}$ and is nonzero for every
adjacent pair $(E,H),(H,Q),(Q,P)$ except $(H,Q)$, treated separately. For a
nonzero slope, $\vhat(s)=\mu+\beta(s-\mu)\to\pm\infty$ as $s\to\pm\infty$, while
$\Ccost(s)=C_0\exp(-\chivar(s-\mu)/\ss)$ is bounded on the half-line where
$\vhat$ dominates: as $s\to+\infty$ the linear term forces $D_{ab}\to+\infty$ for
the higher-slope action; in the $(E,H)$ comparison the cost term is absent, and
in the $(Q,P)$ comparison the common engagement cost cancels, so the slope gap
$\eta>0$ on $\vhat$ governs the crossing, which is finite. The $(H,Q)$ comparison
has $B(H)=B(Q)$ and reduces to $\Ccost(s)=A(Q)-A(H)$; since $\Ccost$ is
continuous, strictly positive, strictly decreasing (Assumption \ref{ass:A2},
$\chivar>0$) with range $(0,\infty)$, this has a unique finite solution whenever
$A(Q)-A(H)>0$, and otherwise Quiet Voice is never optimal and the cutoff sits at
the bracket endpoint. Taking $\underline{s}$ below and $\overline{s}$ above all
(finitely many, uniformly bounded) crossings gives the claim.
\end{proof}

\begin{lemma}[Self-map and ordering preservation]
\label{lem:selfmap}
Under Assumptions \ref{ass:A1}, \ref{ass:A2}, \ref{ass:A4} the map $\Tmap$ is
continuous and maps $\Theta$ into itself, $\Tmap(\Theta)\subseteq\Theta$,
including the collapse faces.
\end{lemma}

\begin{proof}
\emph{Continuity.} The action probabilities $\omega_\bullet$ are continuous in
$\kvec$ (compositions of the smooth signal CDF $G$); the posteriors, prices and
takeover probability are continuous in $(\kvec,\omega)$ on each information cell;
and each $D_{ab}$ in \eqref{eq:diff} crosses zero transversally (strictly
monotone in $s$ at the crossing, by the slope/cost argument of
\ref{lem:bounded}), so the implicit-function theorem gives a solution continuous
in $\kvec$. On a collapse face the relevant $D_{ab}$ is degenerate; there we
define $\Tmap$ by the limiting boundary cutoff, the continuous extension. The
operative $0/0$ object $\pibar=\omegaQ/(\omegaH+\omegaQ)$ extends continuously
through the collapse face because numerator and denominator are Lipschitz and the
denominator is bounded below away from the double-collapse corner, where
l'Hopital gives a direction-independent limit.

\emph{Into the bracket.} By \ref{lem:bounded} every component of $\Tmap(\kvec)$
lies in $[\underline{s},\overline{s}]$.

\emph{Ordering.} We must show
$\Tmap_1(\kvec)\le\Tmap_2(\kvec)\le\Tmap_3(\kvec)$. This is the content of the
global monotone--best-response lemma (\ref{lem:monotone}): the upper envelope
$\max_{(q,a)}U(q,a;s)$ is a piecewise combination of the affine-in-$\vhat$
functions \eqref{eq:payoff-affine} with totally ordered slopes
$B(E)<B(H)=B(Q)<B(P)$, hence the optimal action is nondecreasing in $s$ and the
indifference signals are automatically ordered. On a collapse face the ordering
holds with equality, so $\Tmap$ maps each face into itself or its interior. Thus
$\Tmap(\Theta)\subseteq\Theta$.
\end{proof}

\begin{proposition}[Existence]
\label{prop:existence-new}
Under Assumptions \ref{ass:A1}, \ref{ass:A2}, \ref{ass:A4}, an equilibrium in
cutoff strategies exists.
\end{proposition}

\begin{proof}
$\Theta$ in \eqref{eq:theta} is nonempty, compact and convex, and by
\ref{lem:selfmap} $\Tmap:\Theta\to\Theta$ is continuous. Brouwer's fixed-point
theorem yields $\kvec^\star\in\Theta$ with $\Tmap(\kvec^\star)=\kvec^\star$, an
equilibrium. The fixed point may lie on a collapse face, in which case the
corresponding action carries zero probability; this is precisely the baseline
($\kone=\kzero$, Hold collapses).
\end{proof}

% ---------------------------------------------------------------------
\subsection{Public voice dominates quiet voice in slope}
% ---------------------------------------------------------------------

\begin{lemma}[$B_P>B_Q$]
\label{lem:BP-BQ}
Public Voice is strictly more sensitive to the signal than Quiet Voice,
$B(P)>B(Q)$, whenever the disclosed-flow takeover probability satisfies
$p(X,1)\le\tfrac12$ for the realized disclosed flow $X$; a sufficient primitive
condition is
\begin{equation}
\label{eq:BPBQ-suff}
\Sbar-\Kcost-P(X,1)\;\le\;\mt .
\end{equation}
\end{lemma}

\begin{proof}
The payoff slopes in $\vhat$ differ between Public and Quiet Voice through two
channels. (1) \emph{Trading:} Public Voice carries $q=+1$, contributing $+1$ to
the slope (the price is taken as given by the price-taking blockholder), against
$0$ for Quiet Voice. (2) \emph{Takeover:} both carry the engagement term
$(\mt-m_0)\,p$; under disclosure ($D=1$ for Public) the takeover probability is
$p(X,1)=1-\Phi(T)$ with $T=(\mbar(X,1)+\Kcost-\Sbar+P(X,1))/\sigxi$, and
$\mbar(X,1)$ rises with the disclosed posterior, which rises with $s$. Hence
$\partial p/\partial\vhat=\phi(T)(\partial\mbar/\partial\vhat)/\sigxi\ge0$, so
the takeover channel adds a nonnegative amount to the Public slope. Therefore
\[
\eta=\underbrace{1}_{\text{trading}}
+\underbrace{(\mt-m_0)\,\frac{\phi(T)}{\sigxi}\,
\frac{\partial\mbar}{\partial\vhat}}_{\ge0}\;>\;0,
\]
giving $B(P)=1+\eta>1=B(Q)$ from the trading channel alone. The condition
$p(X,1)\le\tfrac12\iff T\ge0\iff$ \eqref{eq:BPBQ-suff} (using
$\mbar(X,1)=\mt$ on the disclosed branch) additionally places the takeover
channel on the rising part of $\Phi$, making the slope gap robust to
perturbations of the disclosed price.\footnote{The unconditional part
($\eta\ge1$) is all the ordering arguments require; \eqref{eq:BPBQ-suff} is the
transparent primitive bound and holds at the baseline, where the disclosed
takeover probability is $p^\star=0.8473$ ($T=-1.025$ at $P$ normalized to zero;
in-line check below), so the trading channel secures $\eta>0$ regardless.}
\end{proof}

% ---------------------------------------------------------------------
\subsection{Global monotone best response (replaces self-confirming)}
% ---------------------------------------------------------------------

\begin{lemma}[Global monotone best response]
\label{lem:monotone}
Fix any measurable bounded pricing function $P(\cdot,\cdot)$ and any takeover
probability $p(\cdot)\in[0,1]$, hence any intercepts $A(q,a)$ in
\eqref{eq:payoff-affine}. Under Assumptions \ref{ass:A1}, \ref{ass:A2} with the
slope ordering \eqref{eq:slopes}, $B(E)<B(H)=B(Q)<B(P)$, the action
correspondence $a^\star(s)\in\arg\max_{(q,a)}U(q,a;s)$ is nondecreasing in $s$ in
the order $E\prec H\prec Q\prec P$, and the upper envelope
$V(s)=\max_{(q,a)}U(q,a;s)$ is convex in $\vhat(s)$. Consequently the best
response is a cutoff partition with ordered thresholds $\kone\le\kzero\le\kD$.
\end{lemma}

\begin{proof}
Each $U(q,a;\cdot)$ in \eqref{eq:payoff-affine} is affine in $\vhat$ once the
engagement cost is folded in: for non-engagement actions ($E,H$) it is affine in
$\vhat$ with slope $B\in\{0,1\}$; for the two engagement actions ($Q,P$) the
comparison has the cost cancel, leaving an affine-in-$\vhat$ difference with slope
$B(P)-B(Q)=\eta>0$. The upper envelope of a finite family of affine functions of
$\vhat$ is convex and piecewise affine, and its argmax is the action with the
largest slope among those still attainable, which is nondecreasing in $\vhat$,
hence (as $\vhat$ is strictly increasing in $s$, Assumption \ref{ass:A1})
nondecreasing in $s$. With totally ordered slopes $B(E)<B(H)=B(Q)<B(P)$ the
envelope is traversed in the order $E,\{H\text{ or }Q\},P$. The tie
$B(H)=B(Q)$ is broken by the intercepts: $Q$ adds the engagement net benefit
$A(Q)-A(H)$ less $\Ccost(s)$, decreasing in $s$, so among the two equal-slope
actions $H$ is chosen for low $s$ (high cost) and $Q$ for high $s$ (low
cost)---again a nondecreasing switch. Thus the optimal action is nondecreasing in
$s$ and the thresholds are ordered. No conjecture about the \emph{form} of the
partition is used; only the belief-determined intercepts and the primitive slope
ordering enter, on and off the equilibrium path, so the argument is not
self-confirming.
\end{proof}

\begin{remark}[Why the tie $B(H)=B(Q)$ is benign]
\label{rem:tie}
Hold and Quiet Voice have equal slopes and are separated purely by the decreasing
cost schedule $\Ccost$ (Assumption \ref{ass:A2}), so the $H\!\to\!Q$ switch is
monotone; if the engagement net benefit is too small the $Q$ region is empty and
the equilibrium sits on the collapse face $\kzero=\kD$. This is the mechanism
behind the baseline Hold collapse and is fully consistent with
\ref{lem:monotone}.
\end{remark}

% ---------------------------------------------------------------------
\subsection{Dropping A7: conditioning the bidder on the price}
\label{ssub:dropA7}
% ---------------------------------------------------------------------

\begin{lemma}[Price-conditioned bidder; A7 redundant]
\label{lem:dropA7}
Suppose the bidder conditions only on the equilibrium price $P$ (and not
separately on $D$ or the stake). Then the bidder's takeover decision and
probability coincide with those under A7. Hence A7 can be dropped.
\end{lemma}

\begin{proof}
The competitive market maker sets $P(X,D)=\deltaliq\,\E[Y\mid X,D]$. Because the
map $(X,D)\mapsto P$ is injective on the realized support at any nondegenerate
calibration (distinct order-flow/disclosure cells command distinct prices, as
$\E[Y\mid\cdot]$ separates cells whenever their posteriors differ), the price
carries all the information in $(X,D)$ that is payoff-relevant to the bidder. The
bidder bids iff $\Pi_B=\Sbar-P+\xibid-\mbar-\Kcost\ge0$ and, integrating over
$\xibid\sim\mathcal N(0,\sigxi^2)$,
\begin{equation}
\label{eq:pbid-onP}
p(P)=1-\Phi\!\left(\frac{\mbar(P)+\Kcost-\Sbar+P}{\sigxi}\right),
\end{equation}
with $\mbar(P)=\E[m\mid P]$. By the law of iterated expectations
$\E[m\mid P]=\E[\E[m\mid X,D]\mid P]=\E[m\mid X,D]$ on the price-injective
support, so $\mbar(P)=\mbar(X,D)$ and \eqref{eq:pbid-onP} equals the A7
expression. On the measure-zero set where two cells share a price the bidder uses
the price-conditional posterior, the correct Bayesian object, which weakly
dominates any cell-specific rule. Thus the bidder problem and $p$ are unchanged,
and A7 is redundant.
\end{proof}

% ---------------------------------------------------------------------
\subsection{Splitting A5: inner-price and cross-cutoff contraction}
\label{ssub:price}
% ---------------------------------------------------------------------

\begin{assumption}[A5a: inner-price contraction]
\label{ass:A5a}
On each information cell the price map $\Psifp(P)=\deltaliq\,\E[Y\mid X,D;P]$ is a
contraction in $P$:
\begin{equation}
\label{eq:A5a}
\bigl|\Psifp'(P)\bigr|
=\deltaliq\Bigl[\,\underbrace{p^\star}_{\text{(i) direct}}
+\underbrace{\bigl|P+\mbar-\Vhat\bigr|\,\tfrac{\phi(0)}{\sigxi}}_{\text{(ii) feedback}}\,\Bigr]
<1 .
\end{equation}
\end{assumption}

\begin{assumption}[A5b: cross-cutoff stability]
\label{ass:A5b}
The induced prices are Lipschitz in the cutoffs with a modulus small enough that
the composed outer map (cutoffs $\to$ prices $\to$ cutoffs) does not amplify
perturbations beyond the bound used in Assumption \ref{ass:A6}.
\end{assumption}

\begin{remark}[Baseline margins for A5a, stated honestly]
\label{rem:A5margins}
Term (ii) of \eqref{eq:A5a}---the feedback of the price on the bidder's
probability and hence on $\E[Y]$---is bounded at the baseline by $0.9474$
(\texttt{phase0\_robustness.md}, \S4). But the sufficient rule \eqref{eq:A5a}
\emph{also} contains the direct term $\deltaliq\,p^\star$, and at the baseline
disclosed posterior the bidder probability is $p^\star\approx0.847$, so
$\deltaliq\,p^\star\approx0.805$ (in-line check below). Adding the two terms,
$\deltaliq[p^\star+(\text{ii})]\approx0.805+0.947>1$, so the displayed feedback
number does \emph{not} by itself establish $|\Psifp'|<1$. \textbf{We therefore
state Assumption \ref{ass:A5a} as a maintained assumption}, not a derived
inequality: it holds whenever
$\deltaliq[p^\star+|P+\mbar-\Vhat|\phi(0)/\sigxi]<1$, which the baseline does not
satisfy under this conservative sufficient rule. The assumption is nonetheless
mild---the rule bounds $\phi(T)$ by $\phi(0)$ and treats both channels as
additive at their worst case---and the inner fixed point is found numerically
without difficulty at every $\kp$ (residuals $\sim10^{-10}$ at fully converged
points; \ref{rem:numreg}). The reader should read \eqref{eq:A5a} as the
transparent primitive condition and $p^\star=0.847$ as the actual baseline slack
on the direct channel.
\end{remark}

% ---------------------------------------------------------------------
\subsection{Uniqueness as a numerical regularity}
% ---------------------------------------------------------------------

\begin{remark}[Numerical regularity, not a proposition]
\label{rem:numreg}
We do not claim global uniqueness as a theorem. Assumption \ref{ass:A6}
(contraction of $\Tmap$) would deliver it, but we cannot derive
$\|D\Tmap\|_\infty<1$ from primitives without further structure on the
cross-cutoff Jacobian. At the baseline calibration the equilibrium located by the
solver is numerically unique: a multistart search with $30$ random ordered seeds
at $\kappa\in\{0.2,0.5,0.8\}$ converges in all $30/30$ cases to a single fixed
point per $\kappa$ (residuals $\sim4\times10^{-10}$;
\texttt{phase0\_robustness.md}, \S2c), and at the hump peak $\kappa^\star=0.59$
cold-start and both warm-started solves agree to $\Dmin$-spread
$6.9\times10^{-9}$ (\texttt{phase0\_facts.md}, \S2). We label this
\emph{numerical regularity} at the calibration, not a uniqueness theorem; the
comparative statics below are shown to hold across all fixed points the
multistart locates.
\end{remark}

% ---------------------------------------------------------------------
\subsection{Endpoint symmetry of the minority gain}
% ---------------------------------------------------------------------

We now show that the headline object $\Dmin(\kp)=\E[\mR(a)\,\1\{\text{bid}\}]$
takes the \emph{same} value in the two no-leakage limits $\kp\to0$ and
$\kp\to1$. We emphasize at the outset what is and is not claimed: \textbf{we
claim only the endpoint-limit equality
$\lim_{\kp\to0}\Dmin(\kp)=\lim_{\kp\to1}\Dmin(\kp)$, not pointwise reflection
symmetry $\Dmin(\kp)=\Dmin(1-\kp)$.} Pointwise reflection is false---the peak is
at $\kappa^\star\approx0.59\neq\tfrac12$ and $\E[\pi\mid D=0]$ is asymmetric about
$\tfrac12$ ($0.655$ at $\kappa=0.25$ versus $0.612$ at $\kappa=0.75$;
\texttt{phase0\_facts.md}, \S3)---so the proof uses no involution and no
general-$\kp$ commutation.

\begin{lemma}[Endpoint information equivalence]
\label{lem:endpoint-info}
On the nondisclosed branch ($D=0$) the order flow is $X=q+z$ with
$q\in\{-1,0\}$ and $z\in\{-1,0,+1\}$, $\PP(z=0)=1-\kp$, $\PP(z=\pm1)=\kp/2$. In
the limit $\kp\to0$ the noise concentrates at $z=0$; in the limit $\kp\to1$ it
concentrates on the wings $z=\pm1$. In \emph{both} limits the realized $D=0$
posterior law converges to the \emph{same} two-point law $\{0,\pibar\}$ with
$\pibar=\omegaQ/(\omegaH+\omegaQ)$.
\end{lemma}

\begin{proof}
The market maker's posterior on engagement given $X$ is determined by the
conditional law of $q$ given $X$.

\emph{Limit $\kp\to0$.} Here $\PP(z=0)\to1$, so $X\to q$ a.s.: $X=-1$ reveals Exit
($\pi=0$) and $X=0$ reveals Hold-or-Quiet, on which the engagement posterior is
$\pibar$. The realized $D=0$ law is the two-point law $\{0,\pibar\}$ weighted by
the action masses.

\emph{Limit $\kp\to1$.} Here $\PP(z=\pm1)\to\tfrac12$ and $\PP(z=0)\to0$, so each
action is shifted to a wing: $q=-1$ produces $X\in\{-2,0\}$ and $q=0$ produces
$X\in\{-1,+1\}$, each with probability $\tfrac12$. The flow that can come only
from $q=-1$ ($X=-2$) carries $\pi=0$; the flow that can come only from $q=0$
($X=+1$) carries $\pi=\pibar$; the overlapping flow $X=0$ (from $q=-1,z=+1$ or
$q=0,z=0$, the latter of vanishing weight as $\kp\to1$) is dominated by the
$q=-1$ contribution and carries $\pi\to0$. Hence as $\kp\to1$ the realized $D=0$
law again converges to $\{0,\pibar\}$ with the same action-mass weights. The
action masses $\omega_\bullet$ are continuous in $\kp$ and the equilibrium
cutoffs converge to a common ordered triple at both ends---numerically
$(0.1379,0.2596,2.6752)$ at both extremes (\texttt{phase0\_facts.md}, \S1)---so
$\pibar$ and the weights coincide in the two limits.
\end{proof}

\begin{proposition}[Endpoint symmetry of $\Dmin$]
\label{prop:endpoint-symmetry}
Under Assumptions \ref{ass:A1}, \ref{ass:A2}, \ref{ass:A4} and \ref{lem:dropA7},
\begin{equation}
\label{eq:endpoint-eq}
\lim_{\kp\to0}\Dmin(\kp)\;=\;\lim_{\kp\to1}\Dmin(\kp).
\end{equation}
\end{proposition}

\begin{proof}
$\Dmin(\kp)$ is a continuous functional of the realized order-flow law through
(a) the $D=0$ posterior law and (b) the disclosed ($D=1$) branch. By
\ref{lem:endpoint-info} the $D=0$ posterior law converges to the same two-point
law $\{0,\pibar\}$ with the same weights in both limits, so the $D=0$
contribution to $\Dmin$ has the same limit at the two ends. For the disclosed
branch: Public Voice ($q=+1$) triggers $D=1$ deterministically with disclosed
posterior $\mt$ (a $\kp$-constant), so the disclosed contribution is
$(\mt-m_0)\,\omegaP\,p(P_{D=1})$; $\omegaP=1-G(\kD)$ and the disclosed price
depend on $\kp$ only through the equilibrium cutoffs, which converge to the same
ordered triple at both endpoints, so the disclosed contribution also has the same
limit. Adding the two gives \eqref{eq:endpoint-eq}. We do \emph{not} assert a
relabeling involution or a general-$\kp$ commutation: the equality is the common
limit of two convergent sequences, each evaluated at its own endpoint.
\end{proof}

\begin{remark}[Numerical confirmation and the limits caveat]
\label{rem:endpoint-honest}
The exact endpoints are degenerate non-equilibria (residual $\sim10^{-1}$ at
$\kappa\in\{10^{-4},1-10^{-4}\}$; \texttt{phase0\_robustness.md}, \S0), so
\eqref{eq:endpoint-eq} is a \emph{limit} statement, not a solved endpoint value.
On the $1\mathrm{e}{-5}$-gated solver the validly converged near-endpoints
$\Dmin(0.02)=0.066778$ and $\Dmin(0.98)=0.066907$ differ by $1.28\times10^{-4}$
(in-line check below; the small gap reflects the steepness of the curve near the
ends, not a symmetry violation), and the exact-endpoint solves of
\texttt{phase0\_facts.md} (\S1) agree to
$|\Dmin(0^+)-\Dmin(1^-)|=6.86\times10^{-7}$. \textbf{Pointwise symmetry is not
claimed and is in fact false}; the proof of \ref{prop:endpoint-symmetry} is
confined to the limits and contains no general-$\kp$ identity, so it does not
contradict the asymmetric interior.
\end{remark}

% ---------------------------------------------------------------------
\subsection{The conditional hump in \texorpdfstring{$\Dmin(\kp)$}{Delta-min(kappa)}}
% ---------------------------------------------------------------------

Endpoint symmetry forces $\Dmin$ to return to a common value as $\kp$ sweeps from
$0$ to $1$. Whether the interior rises above that common value (a hump) or falls
below it (a trough) is the content of this subsection. We \emph{derive the engine
from the equilibrium objects and verify its sign numerically with the model's own
prices}; in particular, the sign of the relevant Jensen object is the
\emph{opposite} of what earlier drafts claimed, so a careful decomposition is
essential.

On the $D=0$ branch the realized engagement posterior takes finitely many values
in $[0,\pibar]$ (with $\pibar=\omegaQ/(\omegaH+\omegaQ)$; at the baseline Hold
collapses so $\pibar=1$), and the per-state takeover probability is $p(X,0)$.
Write $\E_0[\cdot]$ for the expectation under the realized $D=0$ order-flow law
and $\PP(D{=}0)$ for its total mass. Using $\mbar=m_0+(\mt-m_0)\pi$ and the tower
identity of \ref{app:proof-minority-decomp}, the nondisclosed contribution
decomposes as
\begin{equation}
\label{eq:D0decomp}
\begin{aligned}
\Dmin_{D=0}(\kp)
&=\underbrace{m_0\,\PP(D{=}0)\,\E_0[p]}_{\text{baseline term }\Lam_{0}}
\;+\;\underbrace{(\mt-m_0)\,\PP(D{=}0)\,\E_0[p]\,\E_0[\pi]}_{\text{activism mean term }\Lam_{1}}\\
&\quad+\;\underbrace{(\mt-m_0)\,\PP(D{=}0)\,\Cov_0(\pi,p)}_{\text{realized-weight gap }\Lam_{2}},
\end{aligned}
\end{equation}
where $\Cov_0(\pi,p)=\E_0[\pi p]-\E_0[\pi]\,\E_0[p]$. The disclosed branch
contributes $\Lam_D=(\mt-m_0)\,\omegaP\,p(P_{D=1})$, also $\kp$-dependent through
$\omegaP$ and the cutoffs.

\begin{definition}[Realized-weight gap]
\label{def:Lambda2}
The \emph{realized-weight gap} is the last term of \eqref{eq:D0decomp},
$\Lam_2(\kp)=(\mt-m_0)\,\PP(D{=}0)\,\Cov_0(\pi,p)$. It is the genuine
Jensen/covariance object at the realized weights---\emph{not} the midpoint chord
diagnostic $h(0)-2h(\pibar/2)+h(\pibar)$, which earlier work correctly rejected
because it ignores the realized weights and evaluates the integrand at the
unrealized point $\pibar/2$.
\end{definition}

\begin{remark}[The sign of $\Lam_2$ is negative: the gap fights the hump]
\label{rem:hvsg}
A central correction to earlier drafts: at the baseline, computed with the
model's own equilibrium prices, $\Lam_2<0$ at \emph{every} interior $\kp$ (it
ranges from $-0.0029$ to $-0.0169$ in $\Dmin$ units; in-line check below). The
reason is economic, not a sign slip: higher inferred engagement $\pi$
\emph{deters} bids (\,$\partial p/\partial\pi<0$, the bid-monotonicity property of
the price-decomposition section\,), so $\pi$ and $p$ are \emph{negatively}
correlated across the realized $D=0$ states and $\Cov_0(\pi,p)<0$. Equivalently
$h(\pi)=\pi\,p(\pi)$ is \emph{convex} on the realized support, so its
realized-weight Jensen gap is negative. \textbf{The hump is therefore \emph{not}
driven by concavity of the integrand}; the gap term works against it. This
refutes both the rigor memo's ``concavity of $g=\mbar p$'' story and the
attempt-3 ``$\Lam_2<0$ via U-shaped variance drives the hump'' story---in fact
$\Var_0(\pi)$ is U-shaped (minimal $0.076$ at the peak $\kappa^\star=0.59$), so
posterior dispersion is \emph{lowest} exactly where $\Dmin$ is highest.
\end{remark}

\begin{proposition}[Conditional hump]
\label{prop:hump}
Let
\begin{equation}
\label{eq:Cstar}
\tag{$C^\star$}
\Gamma(\kp)\;\equiv\;\Dmin(\thetastar(\kp))
=m_0\,\PP(D{=}0)\,\E_0[p]+(\mt-m_0)\,\PP(D{=}0)\,\E_0[\pi p]+\Lam_D(\kp)
\end{equation}
denote the full $\kp$-varying value of $\Dmin$ along the equilibrium path
$\thetastar(\kp)$. Suppose $\Gamma$ is single-peaked: strictly increasing on
$(0,\kappa^\star)$ and strictly decreasing on $(\kappa^\star,1)$ for a unique
interior $\kappa^\star$. Then, with the endpoint equality \eqref{eq:endpoint-eq},
$\Dmin(\kp)$ has a single interior maximizer (a hump) at $\kappa^\star$. If
$\Gamma$ is single-troughed, $\Dmin$ has an interior minimizer (a trough). The
orientation is therefore \emph{conditional} on $(C^\star)$, a hypothesis on the
equilibrium path that must be checked at the calibration.
\end{proposition}

\begin{proof}
By \ref{prop:endpoint-symmetry}, $\lim_{\kp\to0}\Dmin=\lim_{\kp\to1}\Dmin$, and
by \ref{app:proof-minority-decomp} with \eqref{eq:D0decomp},
$\Dmin(\kp)=\Gamma(\kp)$ is continuous on $(0,1)$. Under $(C^\star)$, $\Gamma$
rises then falls with a single turning point $\kappa^\star\in(0,1)$; since the two
endpoint limits coincide, $\Gamma$ exceeds that common value throughout $(0,1)$
and attains a unique interior maximum at $\kappa^\star$. The trough case is
symmetric. We do \emph{not} sign $\Gamma'$ from primitives: it nets a
\emph{rising} mean-bid-incidence channel ($m_0\PP(D{=}0)\E_0[p]$ and the activism
mean term) against the \emph{negative} deterrence covariance $\Lam_2<0$
(\ref{rem:hvsg}) and against the unsigned cutoff-shift term below, and the net
sign is calibration-dependent. This is why $(C^\star)$ is a stated hypothesis.

\emph{The unsigned cutoff-shift term.} In general equilibrium the cutoffs move
with $\kp$, contributing $\sum_i(\partial\Dmin/\partial k_i)(d k_i/d\kp)$ to
$d\Dmin/d\kp$. This term is genuinely unsigned: $\Dmin$ is \emph{minority}
welfare, not the blockholder's objective, so the envelope theorem does not zero
it. Because $(C^\star)$ is imposed on the equilibrium-path object
$\Gamma(\kp)=\Dmin(\thetastar(\kp))$, the cutoff motion is already inside the
hypothesis.
\end{proof}

\begin{remark}[Numerical verification of the engine, with the model's own prices]
\label{rem:Lambda2-numerics}
Computed along the equilibrium path with the model's equilibrium prices and bid
probabilities (in-line check below), the components of \eqref{eq:D0decomp} and the
disclosed branch are
\begin{center}
\small
\begin{tabular}{lccccccc}
\toprule
$\kp$ & $0.10$ & $0.25$ & $0.40$ & $0.50$ & $0.59$ & $0.75$ & $0.98$ \\
\midrule
$\E_0[p]$ (mean bid prob)        & $.366$ & $.404$ & $.427$ & $.435$ & $.436$ & $.421$ & $.341$ \\
$\Var_0(\pi)$                    & $.145$ & $.124$ & $.094$ & $.080$ & $.076$ & $.095$ & $.139$ \\
$\Lam_0$ (baseline)              & $.0362$& $.0396$& $.0413$& $.0419$& $.0419$& $.0409$& $.0338$ \\
$10^3\,\Lam_2$ (gap, $<0$)       & $-16.9$& $-14.5$& $-10.8$& $-9.0$ & $-8.4$ & $-10.6$& $-15.9$ \\
$10^3\,\Lam_D$ (disclosed)       & $.014$ & $.044$ & $.101$ & $.137$ & $.146$ & $.087$ & $.008$ \\
$\Dmin$                          & $.0679$& $.0718$& $.0754$& $.0770$& $.0776$& $.0754$& $.0669$ \\
\bottomrule
\end{tabular}
\end{center}
Three facts establish the engine. (i) The mean nondisclosed bid probability
$\E_0[p]$ and the baseline term $\Lam_0=m_0\PP(D{=}0)\E_0[p]$ are both
\emph{hump}-shaped, peaking at $\kappa^\star\approx0.59$; the disclosed
contribution $\Lam_D$ (tracking $\omegaP$, which also peaks at $0.59$) is likewise
hump-shaped but small. Both branches push the same way, and
$\Dmin_{D=0}+\Lam_D=\Dmin$ exactly (in-line check). (ii) The realized-weight gap
$\Lam_2$ is \emph{negative} throughout (\ref{rem:hvsg}); it fights the hump but is
dominated by the rising baseline and activism-mean terms; and $\Var_0(\pi)$ is
U-shaped, so the hump is \emph{not} a posterior-dispersion effect. (iii)
$(C^\star)$ holds at the baseline: $\Gamma$ is single-peaked, so $\Dmin$ humps
with peak $\Dmin(\kappa^\star)=0.07756$ at $\kappa^\star=0.59$ and amplitude
$13.9\%$ of level over the converged grid (in-line check; on the validly-solved
interior of \texttt{phase0\_robustness.md}, \S2a, the amplitude is $15.8\%$). The
orientation is conditional: the $16/20$ hump-vs-trough map of
\texttt{phase0\_robustness.md} (\S2b) finds $4/20$ cells (all at $\sigxi=0.60$)
where $\Gamma$ is single-troughed and $\Dmin$ dips---exactly the trough case of
\ref{prop:hump}, and the reason $(C^\star)$ must remain a stated hypothesis. We do
\emph{not} cite the midpoint-chord column of that map as support, since it is the
rejected diagnostic (\ref{rem:hvsg}).
\end{remark}

\begin{remark}[Across-equilibria robustness]
\label{rem:across-eqm}
The hump is not a branch-selection artifact. At the peak $\kappa^\star=0.59$ the
solver returns the same fixed point from a cold start and from warm starts seeded
from both neighbors, with $\Dmin$-spread $6.9\times10^{-9}$
(\texttt{phase0\_facts.md}, \S2), and warm-starting the unique $\kappa=0.5$ fixed
point across the grid preserves the hump on every examined branch
(\texttt{phase0\_robustness.md}, \S2c). Thus \ref{prop:hump} holds across all
equilibria the multistart locates, consistent with the numerical-regularity
status of uniqueness (\ref{rem:numreg}).
\end{remark}
